The pipeline consists of three components, built and tested independently:
a baseline vision-language model with no retrieved evidence, an image
retriever over the ROCOv2 corpus, and a retrieval-augmented model that
composes the first two. All three share code where possible, which is the
design choice that makes the baseline-vs-RAG comparison in
Section~\ref{sec:results} valid rather than an apples-to-oranges
benchmark.

\subsection{Baseline Vision-Language Model}

The baseline condition uses a pretrained, general-domain vision-language
model checkpoint (a BLIP-2-family model~\cite{li2023blip2}, configured in
\texttt{configs/model\_config.yaml}) with no modification and no retrieved
context: it is given an image and a question, and generates a free-text
answer directly. This is deliberately the simplest possible condition ---
it isolates the effect of retrieval augmentation from any effect of
domain-specific fine-tuning, which prior work has already studied
separately~\cite{li2023llavamed}.

\subsection{Retriever}

The retriever embeds every ROCOv2 image with a CLIP
image encoder~\cite{radford2021clip}, L2-normalizes the resulting vectors,
and indexes them with a FAISS~\cite{johnson2019faiss} flat inner-product
index --- exact (non-approximate) nearest-neighbor search, appropriate at
ROCOv2's corpus scale ($\sim$90K vectors) without requiring an approximate
index such as IVF or HNSW. At query time, the same CLIP encoder embeds the
query image, and the index returns the $k$ most visually similar corpus
entries by cosine similarity (equivalent to inner product over normalized
vectors), together with their captions and similarity scores.

\subsection{Retrieval-Augmented Model}

The retrieval-augmented condition composes the baseline model and the
retriever rather than reimplementing generation: it retrieves the top-$k$
similar corpus captions for the query image, constructs a prompt that
prepends those captions (truncated to a fixed word budget) before the
original question, and passes that prompt through the \emph{same}
underlying generation call the baseline condition uses. Concretely, both
conditions call one shared code path
(\texttt{BaselineVLM.generate\_from\_prompt} in the released code) with the
same tokenization, decoding parameters (deterministic decoding; see
Section~\ref{sec:experimental_setup}), and post-processing --- the only
difference between the two experimental conditions is the text of the
prompt itself. If no evidence is retrieved for a given query (e.g.\ a
cold or empty index), the system falls back to the bare question with no
evidence block, i.e.\ to baseline behavior, rather than emitting a
malformed prompt.

\subsection{Design Rationale: Shared Generation Code}

This shared-code-path design is the paper's central methodological
safeguard. A comparison between two differently-implemented pipelines ---
even ones intended to differ only in ``using retrieval or not'' --- risks
confounding the retrieval effect with incidental implementation
differences (e.g.\ different tokenization, different truncation behavior,
different decoding temperature). By construction, that confound is
eliminated here: the two conditions differ in exactly one respect, the
prompt text, which is what licenses attributing any measured accuracy
difference (Section~\ref{sec:results}) to retrieval augmentation itself.
